\section{黑水城文书让谁进入西夏史}

1049--1193年的连续在位纪年覆盖李谅祚、李秉常、李乾顺、李仁孝与李纯佑等时期。宋、辽、金史能够为王位更替和外交事件提供外部年序，却很少让我们看见西夏境内的法律、寺院、印刷、债务与基层行政。黑水城及河西、宁夏的西夏文、汉文、藏文与回鹘文文书，正因此具有主体材料的意义：它们不只是替宋方叙事添几个细节，而改变了可以追问的问题。

不过，文书并不天然代表“西夏社会”。每件材料都须先辨语言、文种、纪年、出土层位或收集史、释文版本和政权归属。黑水城的保存密度很高，不能代替河套、都城或西夏全境；一部法律或一纸官文书说明规范和机构术语，也不自动证明它在每处被照样执行。把这些限制写清，才避免用罕见材料再造一个过于整齐的国家。

\section{朝廷、寺院与边地资源}

在李乾顺至李仁孝的长时期内，西夏与金的册封、贡赐和边境安排构成相对稳定的外部条件。稳定并不等于停止征发：城寨、河西道路、水源、军队与佛寺都需要资源和人员。西夏文法律、行政文书与佛经印本可使我们把“国家”拆解为多种机构的活动；宋金正史、碑刻和考古则可帮助核对地点与外部关系。任何一个来源都不足以独自量化税额、户口或兵额。

宗教也不应被视为与政治无关的装饰。印本、经卷、造像与寺院材料能够显示书写、捐施和跨语言接触，却不让我们直接知道每一位居民的信仰。对党项、汉人或其他史料分类同样如此：官号、姓名、文字和职业分别记录不同事实，不能被合并为固定的现代族群。西夏在1227年覆亡之前的国家能力，正存在于这些不完全重合的文书、城址与边境关系之中。
